\chapter{Bonnes pratiques}
\label{ch:bestpractices}

Leçons concentrées tirées de déploiements réels. Chaque section est
optionnelle — prenez ce qui colle à votre échelle.

\section{Hiérarchie de stockage}

\begin{figure}[ht]
  \centering
  \includegraphics[width=0.95\textwidth]{04-locations-desktop.png}
  \caption{La page \emph{Emplacements} sur une installation
    réelle~: une hiérarchie domestique avec un bureau, une
    bibliothèque dans la pièce, dix étagères et deux pièces de
    plain-pied \string— chaque n\oe ud porte son code L et un
    compteur de volumes à jour.}
  \label{fig:locations}
\end{figure}

\begin{itemize}
  \item \textbf{Trois niveaux suffisent en général.} Pièce →
        étagère → rangée couvre la plupart des foyers. Un quatrième
        niveau (par exemple « boîte à l'intérieur de la rangée »)
        devient de la paperasse à maintenir.
  \item \textbf{Nommez les lieux comme vous en parlez.} « Étagère
        du salon, en haut » bat « SAL-E01-R01 ». Les noms
        apparaissent sur les étiquettes ; traitez-les comme du
        contenu visible.
  \item \textbf{Imprimez les codes L une seule fois.} Réimprimer
        parce que vous avez rééquilibré les étagères est du
        travail perdu — les codes L sont des identifiants
        stables ; le nom humain sur l'étiquette peut dériver et
        c'est OK.
\end{itemize}

\section{Planches de codes V}

\begin{itemize}
  \item Imprimez les codes V en planches de 50 à 100. Des
        autocollants pré-imprimés dans un porte-étiquettes près du
        poste de scan rendent le catalogage \emph{rapide}.
  \item \textbf{Collez le code V avant de scanner.} Cela élimine
        un futur mode de panne où vous mal-étiquetez un livre et
        devez le re-trouver par scan.
  \item \textbf{Placez l'étiquette toujours au même endroit.}
        Intérieur de la 4\textsuperscript{e} de couverture, coin
        supérieur droit, parallèle au dos — choisissez une règle
        et tenez-vous-y. Le futur-vous lors d'un audit remerciera
        le présent-vous.
\end{itemize}

\section{Genre versus Dewey}

mybibli supporte à la fois des genres larges et des codes
Dewey\index{Dewey}. Choisissez l'un des deux comme aide principale
à la navigation :

\begin{itemize}
  \item \textbf{Utilisez les genres} quand le public navigue par
        humeur (« je veux un thriller ») et que la collection est
        assez petite pour que « Thriller » soit suffisamment
        granulaire.
  \item \textbf{Utilisez Dewey} quand le public navigue par sujet
        (« livres sur l'histoire de France du XVIII\textsuperscript{e}
        siècle ») ou quand la collection a atteint la taille où
        « Documentaire » n'est plus un filtre utile.
  \item Utiliser les deux est OK ; mybibli ne force pas le choix.
\end{itemize}

Les deux ne se comportent pas de la même façon lors d'une récupération
de métadonnées, et la différence s'imagine facilement à l'envers :

\begin{itemize}
  \item \textbf{Le genre n'appartient qu'à vous.} Aucune bibliothèque
        de la chaîne n'en fournit, et aucune récupération de
        métadonnées \string— ni celle du scan, ni
        \emph{Récupérer les métadonnées}, ni le rattrapage de
        l'administration \string— ne le modifie. Attribuez-le une fois,
        il reste.
  \item \textbf{Le code Dewey, lui, peut vous être proposé.} Une
        bibliothèque qui détient le titre peut le fournir, et il
        remplira un champ vide. Dès que vous le corrigez vous-même,
        mybibli enregistre que la valeur est la vôtre et cesse de
        l'écraser.
\end{itemize}

Les deux champs sont réunis dans la rubrique \emph{Classement} du
formulaire d'édition des métadonnées, à l'écart des champs
bibliographiques que décrivent les bibliothèques, et cette rubrique
rappelle la distinction à l'écran.

\section{Cadence d'audit}

\begin{itemize}
  \item \textbf{Trimestriel} pour les collections sous quelques
        centaines de volumes.
  \item \textbf{Mensuel} pour les collections au-dessus,
        spécialement si plus d'une personne catalogue.
  \item Lancez toujours un audit \emph{après} un grand
        rééquilibrage de rayons — déplacer des livres
        physiquement sans le dire à mybibli est la première
        source de bugs « où est ce volume ? ».
\end{itemize}

\section{Comptes utilisateur}

\begin{itemize}
  \item \textbf{Un humain, un compte.} Partager un seul compte
        \texttt{admin} entre deux personnes fait perdre la piste
        d'audit et rend la garde du dernier admin moins utile.
  \item \textbf{La plupart des éditeurs doivent être
        Bibliothécaires, pas Admins.} Le panneau admin gère des
        préoccupations d'installation et de nettoyage rare ; vous
        n'avez pas besoin de droits admin pour cataloguer ou
        prêter.
  \item \textbf{Changez le mot de passe admin seedé au premier
        login.} Sur 1.1.0+ l'assistant le force ; sur les
        installations 1.0.x existantes, faites-le manuellement
        avant d'exposer l'instance sur un réseau.
\end{itemize}

\section{Prêts}

\begin{itemize}
  \item \textbf{Gardez des emprunteurs identifiables.} « Marie L. »
        bat « Marie » dès la seconde Marie qui pointe son nez.
        Email ou téléphone dans le champ notes aide pour les
        relances.
  \item \textbf{Parcourez la liste des retards une fois par
        semaine.} L'indicateur du dashboard vous houspille pour
        une raison — les prêts en retard cumulés sur un an
        deviennent des livres que vous ne possédez plus.
  \item \textbf{Notez les dégâts au retour.} Basculer un volume
        retourné en \emph{Usé} ou \emph{Endommagé} capture
        l'information pendant que l'emprunteur est devant vous ;
        chasser le coût plus tard de mémoire fonctionne rarement.
\end{itemize}

\section{Fournisseurs de métadonnées}

\begin{itemize}
  \item \textbf{Limitez les clés.} Vous n'avez pas besoin des
        trois fournisseurs à clé dès le premier jour — BnF + Open
        Library couvrent les livres gratuitement. Ajoutez OMDb +
        TMDb le jour où vous cataloguez votre premier DVD.
  \item \textbf{Traitez les clés comme des secrets.} \texttt{.env}
        est dans \texttt{.gitignore} ; conservez les sauvegardes
        de \texttt{.env} chiffrées (chapitre~\ref{ch:backup}).
\end{itemize}

\section{Mises à jour}

\begin{itemize}
  \item \textbf{Lisez les notes de version avant de tirer
        \texttt{:latest}.} Voir chapitre~\ref{ch:upgrade}.
  \item \textbf{Sauvegardez juste avant chaque mise à jour.} Même
        les migrations « additives » ont une surface de panne
        non-nulle ; une sauvegarde minute-avant transforme une
        inquiétude en formalité.
  \item \textbf{Épinglez une version pour la production.} Les
        auto-updates conviennent à une instance familiale ; pour
        une bibliothèque communautaire, un saut majeur non
        surveillé peut interrompre le service.
\end{itemize}
